\section{讨论与局限}\label{sec:discussion}

\textbf{该层带来了什么。} 与其数量无关代价的硬约束、在秩亏与弱活跃接触下仍存活的导数、以及与 $p$ 无关代价的全局性后验证据——物理可容许性、可训练性与失效信号。\textbf{它没有带来什么}：前向求解的全局性。ALM 是 KKT 方法，证书可能失败。我们认为这是正确的取舍，因为一个会说“未认证”的层，比一个会错误地说“已认证”的层更有用。

\textbf{四腿例子比看上去更窄。} 旗标性的 LICQ 失效案例在参数使支撑面刚性移动时可微（且我们的方法是精确的），而在接触高度相互独立时不可微。在我们的真实数据协议中，支撑平面每帧由深度估计一次，因此是\emph{刚性}输入参数，相容性成立；非刚性情形出现在四个接触高度可学习时，而我们不这样做。

\textbf{证书的偏差与方差。} 松弛在我们的合成实例与无约束子情形上是紧的，但紧性是经验的，不是定理。在真实数据上，它以依赖数据集的比率变紧：基础割下为 24.4/15.3/13.1\%，可靠 bound-RLT 族下为 51.1/41.7/23.8\%（表~\ref{tab:rigorous}）。在失败处，主张退化为“可靠下界加可行性检查”，它仍能界定次优性，但不是认证；未认证行是真正的松弛（中位间隙 $1$--$2\times10^{-5}$），而非数值上无法定论。可靠性不再是经验的：借助定理~\ref{thm:cert} 的对偶修复，每个报告的严格间隙按构造非负，因此失效模式是\emph{松弛}，绝不是错误证书。

\textbf{范围。} 非仿射约束（距离、间隙）会破坏常维约简，因为 $g_j^2$ 不再是 $(\mathbf v_2,\mathbf t)$ 中的二次型；分支交叉情形需要混合/Laplace 模型，那是本计划中的另一篇论文。该方法刻意与另外两个原型正交：它假设单个受约束最优解与单峰后验，并把预算花在约束与秩鲁棒微分上。
